\medskip
\noindent
\SGARefTarget{sga7.sga7i.eix.n12-6}{12.6}. Notice that this statement proves \SGARefLink{sga7.sga7i.eix.n10-4}{10.4} in the particular case
under consideration when, in \SGARefLink{sga7.sga7i.eix.n10-4}{10.4} b), one takes for $\xi$ the canonical
polarization considered there. As we showed in \SGARefLink{sga7.sga7i.eix.n10-5}{10.5}, this suffices to
imply \SGARefLink{sga7.sga7i.eix.n10-4}{10.4} in the general case.

\medskip
\noindent
\SGARefTarget{sga7.sga7i.eix.n12-7}{12.7}. \underline{Proof of \SGARefLink{sga7.sga7i.eix.n12-5}{12.5}. The case $\ell\ne p$.} It relies
essentially on the \underline{cohomological Picard--Lefschetz formula}
in the situation considered in \SGARefLink{sga7.sga7i.eix.n12-5}{12.5}, a formula giving the action of the
monodromy group $I$ on the following group:\footnote{The working
transcription adds a dual superscript to $T_\ell(A_{\bar\eta})$ in
(\SGARefLink{sga7.sga7i.eix.n12-7-1}{12.7.1}). A direct 5,000-dpi inspection of the printed formula shows no
such superscript, in agreement with the standard identification
$H^1(A_{\bar\eta},\mathbb Z_\ell)=
\operatorname{Hom}(T_\ell(A_{\bar\eta}),\mathbb Z_\ell)$.}

\begin{equation}\tag{\SGARefTarget{sga7.sga7i.eix.n12-7-1}{12.7.1}}
 H^1(X_{\bar\eta},\mathbb Z_\ell)
 \simeq H^1(A_{\bar\eta},\mathbb Z_\ell)
 \simeq\operatorname{Hom}
   \bigl(T_\ell(A_{\bar\eta}),\mathbb Z_\ell\bigr)
\end{equation}

\noindent
when $\ell\ne p$. This formula, which will be proved in a later expos\'e
(Expos\'e XV), and indeed in a substantially more general case, is
essentially identical to assertion \SGARefLink{sga7.sga7i.eix.n12-5}{12.5} for the prime $\ell$ in
question, in view of the definition in \SGARefLink{sga7.sga7i.eix.n9-2}{9.2} of $u_\ell$ for $\ell\ne p$.
It therefore remains to treat the case $\ell=p$. We shall do this by a
method which reduces it to characteristic zero (and, moreover, to the
case where $X_o$ has only one singular point), which we now explain.

\medskip
\noindent
\underline{Remarks \SGARefTarget{sga7.sga7i.eix.n12-7-2}{12.7.2}.} Notice also that, once the announced
reduction has been made, the usual arguments would allow us to reduce to
the case in which the base field is the field $\mathbb C$ of complex
numbers, and then to invoke the Picard--Lefschetz formula in its classical
form---except that there does not seem to be a satisfactory reference
for that formula. The proof of the Picard--Lefschetz formula given in our
Seminar will proceed similarly, by reduction to the transcendental case,
which is amenable to properly transcendental methods. One may therefore
say that, in any event, the proof of \SGARefLink{sga7.sga7i.eix.n12-5}{12.5} (and hence of \SGARefLink{sga7.sga7i.eix.n10-4}{10.4}) proposed
here is transcendental in nature. Let us also point out that the only
other known proof of \SGARefLink{sga7.sga7i.eix.n10-4}{10.4} comes from rigid-analytic theory [24], and must
for that reason likewise be regarded as essentially ``transcendental.''
It would be interesting to find purely algebraic proofs of these results
(although it is by no means evident that this is possible!).
